\chapter{CONCLUSIONES, RECOMENDACIONES Y BIBLIOGRAFIA}

\section{Conclusiones}

\begin{enumerate}
  \item La vivienda unifamiliar de tres pisos ubicada en Capachica se atiende mediante una instalación eléctrica residencial monofasica de 220 V, con una sectorización ordenada de 6 circuitos exclusivos (alumbrado y tomacorrientes por nivel).
  \item La demanda máxima estimada del proyecto alcanza los 4.30 kW, con una corriente demandada de 21.73 A, lo que justifica la instalación de un interruptor general termomagnético de 2P-32A y un interruptor diferencial de seguridad general de 2P-40A (30mA).
  \item Se ha estructurado el sistema con un Tablero General TG-01 en el primer piso y dos tableros secundarios de distribución TD-01 y TD-02 en los niveles superiores, lo que optimiza las longitudes de cableado y facilita la maniobra.
  \item Se han excluido todas las cargas especiales ajenas al levantamiento de puntos (lavandería, motor o bomba de agua), logrando un diseño simplificado e integrado para los tres pisos.
  \item La selección de conductores (1.5 mm2 para alumbrado, 2.5 mm2 para tomacorrientes y 10 mm2 para la acometida alimentadora) cumple satisfactoriamente con la capacidad de corriente y reserva normada.
  \item La incorporación de interruptores diferenciales en circuitos de tomacorrientes y la conexión franca a un pozo de puesta a tierra con resistencia inferior a 15 Ohms aseguran la protección eficaz de los usuarios.
\end{enumerate}

\section{Recomendaciones}

\begin{enumerate}
  \item Validar en campo la ubicación final y el recorrido físico de las canalizaciones en losa y piso antes del vaciado de concreto.
  \item Confirmar con la empresa distribuidora de energía de la zona la ubicación del medidor y la acometida en fachada.
  \item Realizar la medición física de la resistencia del pozo de puesta a tierra tras su construcción para garantizar un valor inferior a 15 Ohms.
  \item Mantener rotulado de forma permanente y clara el tablero general TG-01 y los tableros secundarios de distribución.
  \item No sobrepasar el número de conductores por ducto según los factores de llenado del CNE-U.
\end{enumerate}

\section{Bibliografia}

\begin{enumerate}
  \item Ministerio de Energia y Minas. \textit{Codigo Nacional de Electricidad -- Utilizacion}. Resolucion Ministerial N. 0037-2006-MEM.
  \item Plataforma del Estado Peruano. \textit{Codigo Nacional de Electricidad -- Utilizacion, PDF oficial}.
  \item Ministerio de Vivienda, Construccion y Saneamiento. \textit{Reglamento Nacional de Edificaciones -- RNE}.
  \item Plataforma del Estado Peruano. \textit{EM.010 Instalaciones Eléctricas Interiores}.
  \item INACAL. \textit{Normas Tecnicas Peruanas (NTP)}.
  \item INACAL. \textit{Listado de normas tecnicas citadas en reglamentos obligatorios}.
  \item INACAL. \textit{INACAL y la estandarizacion}.
\end{enumerate}

\bigskip

\noindent\textbf{Nota final:} Este documento tiene fines exclusivamente académicos y no constituye un proyecto ejecutivo para construccion. Su contenido debe ser revisado y ajustado por un profesional habilitado antes de cualquier implementacion real.
